\documentclass[aspectratio=169]{beamer}

% ── Theme ──────────────────────────────────────────────────────────────────
\usetheme{metropolis}
\usepackage{appendixnumberbeamer}
\usepackage{booktabs}
\usepackage{hyperref}

% ── Colours ────────────────────────────────────────────────────────────────
\definecolor{RIVMblue}{HTML}{003D6E}
\definecolor{RIVMaccent}{HTML}{0070BB}
\definecolor{lightgray}{HTML}{F5F5F5}

\setbeamercolor{frametitle}{bg=RIVMblue, fg=white}
\setbeamercolor{progress bar}{fg=RIVMaccent}
\setbeamercolor{title separator}{fg=RIVMaccent}
\setbeamercolor{alerted text}{fg=RIVMaccent}

% ── Metadata ───────────────────────────────────────────────────────────────
\title{Artificial Intelligence and Public Health}
\subtitle{The changing landscape of infectious disease surveillance}
\author{Albert Jan van Hoek}
\institute{RIVM --- National Institute for Public Health and the Environment}
\date{23 March 2026, Amsterdam}

% ══════════════════════════════════════════════════════════════════════════
\begin{document}

\maketitle

% ── Slide 2 ────────────────────────────────────────────────────────────────
\begin{frame}{Who am I?}
  \begin{itemize}
    \item Senior researcher at RIVM (Dutch National Institute for Public Health)
    \item PI of \textbf{Infectieradar} --- a large online cohort reporting weekly
          symptoms of acute infections, with virological testing of $\sim$200
          participants per week
    \item Infectious disease modelling and health economics
    \item Software architecture and functionality (no coding)
    \item Experimenting with AI: prompting, NLP, agentic AI --- trying to
          understand and make sense of it
  \end{itemize}
\end{frame}

% ── Slide 3 ────────────────────────────────────────────────────────────────
\begin{frame}{Prompting AI}
  \begin{center}
    \large Writing a good prompt is writing a good plan.
  \end{center}
  \vspace{1em}
  \textbf{The key: level of abstraction.}
  \vspace{1em}

  \textit{Example: ``come to dinner'' in Dutch}

  \medskip
  \textbf{Literal prompt:} ``Translate `come to dinner' into Dutch''

  \medskip
  \textbf{Abstract prompt:}
  \begin{quote}
  ``Write a short call to signal that food is ready and people can come to
    the table. Keep it concise, express mild urgency, use language level B2,
    in Dutch.''
  \end{quote}

  The second version has \emph{adjustable parameters}: urgency, register,
  language level, language. Swap one, get a different result --- in French,
  Turkish, Chinese, or Dutch.
\end{frame}

% ── Slide 4 ────────────────────────────────────────────────────────────────
\begin{frame}{What is AI, at the right level of abstraction?}
  \begin{itemize}
    \item A process that answers natural-language questions using a trained
          neural network
    \item Trained on large amounts of human-generated data with human feedback
          (Reinforcement Learning from Human Feedback --- RLHF)
    \item Our collective knowledge is captured in the model
    \item The knowledge lives in hardware; electricity running through a trained
          network produces the answer
    \item \textbf{No magic. Silicon and electricity.}
  \end{itemize}
  \vspace{0.5em}
  \alert{Training is unpredictable.} Capacity improves predictably in aggregate,
  but what exactly will emerge at the next scale cannot be specified in advance.
\end{frame}

% ── Slide 5 ────────────────────────────────────────────────────────────────
\begin{frame}{Reinforcement learning --- machines}
  \textbf{Training a neural network:}
  \begin{enumerate}
    \item Start with an untrained network
    \item Present data \hfill \textit{(pictures: cats / dogs)}
    \item Adjust the network \hfill \textit{(slightly different weights)}
    \item Check function \hfill \textit{(can it distinguish cats from dogs?)}
    \item Feedback \hfill \textit{(metric: how well does it do?)}
    \item Select the more trained network
    \item Present data again / new data
    \item Repeat
  \end{enumerate}
\end{frame}

% ── Slide 6 ────────────────────────────────────────────────────────────────
\begin{frame}{Reinforcement learning --- humans}
  \textbf{Training a human neural network:}
  \begin{enumerate}
    \item Start with an untrained network
    \item Present data
    \item Change in the network
    \item Check function
    \item Feedback
    \item Select the more trained network
    \item Present data again / new data
    \item Repeat
  \end{enumerate}
  \vspace{0.5em}
  As a human you experience this \emph{from the inside}. As a scientist I
  recognise this directly: building a health-economic model is an iterative loop
  of construction, questioning, and revision.
\end{frame}

% ── Slide 7 ────────────────────────────────────────────────────────────────
\begin{frame}{A network mental model}
  In infectious disease modelling I use network models to predict measles
  transmission in school--household networks.

  \begin{itemize}
    \item \textbf{Nodes}: schools (or individuals)
    \item \textbf{Edges}: connections (children sharing a household)
    \item Dots and lines
  \end{itemize}

  \vspace{0.5em}
  Training a neural network looks the same: neurons are dots, connections are
  lines, and with each training round the lines get denser and reconfigure.

  \vspace{0.5em}
  Networks have rules and dynamics. \textbf{The point of view from each node is
  different.} I can imagine myself inside that network --- look around you.
\end{frame}

% ── Slide 8 ────────────────────────────────────────────────────────────────
\begin{frame}{Learning as a network process}
  \begin{itemize}
    \item Learning creates new links; thinking \emph{uses} those links
    \item Thinking is a verb --- a decentralised action across thousands of nodes
    \item ``Passing things through the population'' \emph{is} the thinking
    \item New connections combine and produce something that was not there before
          --- \textbf{emergence}
  \end{itemize}

  \vspace{0.5em}
  From this perspective, AI training is not a unique process. Improvement in AI
  looks like evolution in our own networks. And both have feedback.
\end{frame}

% ── Slide 9 ────────────────────────────────────────────────────────────────
\begin{frame}{Evolution as a learning process}
  Viruses also ``learn''.

  \begin{itemize}
    \item They evolve through variation and selection
    \item That evolution is precisely what keeps them around
    \item What evolves is an emergent process between virus and hosts
    \item That it will evolve is certain; what exactly changes is not ---
          it depends on too many things
  \end{itemize}

  \vspace{0.5em}
  \alert{Same loop. Different substrate. Same unpredictability of outcome.}
\end{frame}

% ── Slide 10 ───────────────────────────────────────────────────────────────
\begin{frame}{What this implies}
  This creates a strange feeling about the relationship between the thinking
  I experience and the underlying process.

  \vspace{0.5em}
  \textbf{Something about what I am.}

  \vspace{0.5em}
  AI training is a process we have extensive experience with --- in biology,
  in epidemiology, in our own development as scientists.

  \vspace{0.5em}
  That familiarity is useful. It means we are not starting from zero.
\end{frame}

% ── Slide 11 ───────────────────────────────────────────────────────────────
\begin{frame}{Why this matters}
  \textbf{AI is not a thing. It is a learning process.}

  \begin{itemize}
    \item It keeps improving --- it is not a normal tool
    \item It is a network phenomenon; its capacities are distributed
    \item Feedback is central to the learning
    \item The loop runs again and again, but the connections change each time,
          producing new capacities and functions
    \item \textbf{Prompting} should connect to the right capacities within
          that network
  \end{itemize}
\end{frame}

% ── Slide 12 ───────────────────────────────────────────────────────────────
\begin{frame}{This framework scales}
  The same image --- learning loop, network, emergence, feedback --- applies to:

  \vspace{0.5em}
  \begin{columns}[t]
    \column{0.5\textwidth}
      \begin{itemize}
        \item AI development
        \item Infectious disease surveillance
        \item Evolution
        \item The immune response
      \end{itemize}
    \column{0.5\textwidth}
      \begin{itemize}
        \item Society and institutions
        \item The economy
        \item Our own thinking
        \item Science itself
      \end{itemize}
  \end{columns}

  \vspace{0.5em}
  It is not a metaphor. It is the same underlying process running on
  different substrates.
\end{frame}

% ── Slide 13 ───────────────────────────────────────────────────────────────
\begin{frame}{Back to public health: surveillance and pandemic preparedness}
  \textbf{Where AI enters our work:}
  \begin{itemize}
    \item Natural Language Processing (NLP)
    \item Generative AI
    \item Benchmarking: formal feedback and training loops
    \item Quality assurance of a learning system
  \end{itemize}

  \vspace{0.5em}
  \textbf{Where humans remain essential:}
  \begin{itemize}
    \item Communication and interpretation
    \item Open data
    \item FAIR principles (Findable, Accessible, Interoperable, Reusable)
  \end{itemize}
\end{frame}

% ── Slide 14 ───────────────────────────────────────────────────────────────
\begin{frame}{Pandemic preparedness as a learning system}
  \begin{itemize}
    \item Learning must be \textbf{continuous} --- preparedness degrades
          without it
    \item Quality monitoring is not optional
    \item The problem is multidimensional: biological, logistical, political,
          social
    \item Nature already solved a version of this: \textbf{the immune response}
          is a heavily regulated, multi-layered learning system
    \item Formal concept: \textbf{$k$-cover regulation} --- maintaining
          sufficient distributed monitoring capacity under uncertainty
  \end{itemize}
\end{frame}

% ── Slide 15 ───────────────────────────────────────────────────────────────
\begin{frame}{Your future role: Guardian of the learning process}
  If AI is a learning process --- not a tool --- the human role is not
  operation. It is \textbf{stewardship}.

  \vspace{0.5em}
  \textbf{Technical requirements:}
  \begin{itemize}
    \item Understanding of the underlying process
    \item Abstract thinking and architecture
    \item Domain knowledge that AI cannot self-verify
  \end{itemize}

  \vspace{0.5em}
  \textbf{Personal qualities:}
  \begin{itemize}
    \item No bias in; no bias out
    \item Honesty --- the loop breaks when the signal is corrupted
    \item Error correction and the patience to keep correcting
    \item The capacity to forgive mistakes --- in others and in the system
  \end{itemize}
\end{frame}

% ── Slide 16 ───────────────────────────────────────────────────────────────
\begin{frame}{Conclusion}
  \begin{itemize}
    \item AI is not a tool; it is a learning process
    \item Use it --- but keep focus on the loop, not the output
    \item We have a great deal of prior experience with this kind of process
    \item \textbf{Guardian of the learning process} may be the right frame for
          the role public health professionals will play
    \item I find these interesting times. I am positive.
  \end{itemize}
\end{frame}

% ── Closing ────────────────────────────────────────────────────────────────
\begin{frame}[standout]
  Keep alive what keeps you alive,\\
  and improve it for the next generation.

  \vspace{1.5em}
  \textit{Disco ergo sum.}
\end{frame}

\end{document}